% Actual level-based repair: no uniform coefficient-field bound is used.
% The modular representation and level estimate are proved in the companion
% height-level-repair fragment. This file proves its height consequence.
\paragraph{Transfer from the level, allowing unbounded coefficient degree.}
Let $A/L$ be the non-CM quaternionic surface. Write $R\geq1$ for the
squarefree product of the prescribed rational bad primes, including $2$,
and assume $[L:\mathbb Q]\leq576$. The modularity construction provides
a weight-two newform of level $N$ and a field $L'$ over which the Tate
modules of $A$ and $J=J_1(N)$ share an irreducible constituent. In the
notation of the sign-splitting construction, the auxiliary Dirichlet
character has order
\[
 m\leq 2\prod_{\substack{p\mid R\\p\ne2}}(p-1)\leq2R.
\]
The sign-splitting cochain has values in $\mu_{2m}$. If one instead
uses the full scalar splitting map, it also includes the square root
of the degree map; that factor equals one on $G_L$. Thus in either
notation the restriction to $G_L$ is a character valued in
$\mu_{2m}$, and its kernel field satisfies
\[
 d'=[L':\mathbb Q]\leq2m[L:\mathbb Q]\leq2304R.
\tag{LT1}
\]
The quadratic twist used to remove extra ramification does not increase
this bound, since $2$ divides $2m$. There is no bound here on $m$
independent of $R$.

Faltings' Hom theorem, followed by scalar extension in the Hom space,
gives a nonzero $L'$-homomorphism $A\to J$. Absolute simplicity of $A$
and Poincar\'e reducibility therefore give an $L'$-isogeny
$A\times Z\to J$ for an abelian variety $Z/L'$. Put $g=\dim J$;
then $g\geq2$. The modular genus estimate
\cite[(5.9), p.~13]{vonKaenel2021} gives $g\leq N^2/24$.
The standard Belyi map has degree at most $N^2$ and
\cite[Theorem~1.1.1 and \S2.3]{javanpeykar} gives
\[
 h_F(J)\leq13000000N^{12}.
\tag{LT2}
\]
All heights in this argument use the original stable Faltings metric.

Suppose the companion local argument has proved $N\leq C R^2$,
where $C\geq1$. Then
\[
 g\leq C^2R^4/24,\qquad
 h_F(J)\leq13000000C^{12}R^{24}.
\]
Set $l=\log C$ and $r=\log R$. The published
Gaudron--R\'emond Theorem~1.9(1), M\'emoires SMF 176 (2023),
pp.~7--8, combined with half-logarithmic isogeny-height variation,
gives
\[
 h_F(A)+h_F(Z)\leq h_F(J)+V,
 \quad
 V=2g^2\left(8g^2\log(7g)+\log d'
       +\log\max(1,\log d',h_F(J))\right).
\tag{LT3}
\]
Here the source variety for the isogeny estimate is $J_{L'}$.
The three logarithms in (LT3) satisfy respectively
\[
 \log(7g)\leq2l+4r,\quad
 \log d'<8+r,\quad
 \log\max(1,\log d',h_F(J))<17+12l+24r.
\]
For the last assertion, the bound in (LT2) also dominates $1$ and
$\log d'$; use $8+r\leq9R^{24}$.
It follows that
\[
 V\leq
 \frac{C^8R^{16}}{24^4}(32l+64r)
 +\frac{C^4R^8}{288}(25+12l+25r).
\tag{LT4}
\]
Using $l\leq C^4/4$, $r\leq R^8/8$, and also
$r\leq R^{16}/16$ in the second summand gives
\[
 \frac{V}{C^{12}R^{24}}
 \leq\frac{16}{24^4}+\frac{25+3+25/16}{288}<1.
\tag{LT5}
\]
The Bost lower bound in \cite[(2.2), p.~5]{vonKaenel2021}
gives $h_F(Z)\geq-\tfrac12\dim Z\log(2\pi^2)>-\tfrac32g$.
Consequently
\[
 h_F(A)<(13000000+1+1)C^{12}R^{24}
       <2^{24}C^{12}R^{24}.
\tag{LT6}
\]
This calculation needs no uniform bound on the degree of the newform's
coefficient field and no assertion that all its geometric factors are
isogenous to $A$.

For the companion conductor estimate one may take $C=210^{4000}$.
Since $210<2^8$ and $4608>2^{12}$,
\[
 2^{24}C^{12}=2^{24}210^{48000}
 <2^{384024}<2^{1990656}<4608^{165888}.
\tag{LT7}
\]
Thus the non-CM height inequality is strictly within the original
coefficient budget, once the stated modularity and level estimate have
been established. In the five-contact application $R=B_*M(s)$, so
the radical exponent is exactly $24$, as required.
